% Fuente única de clase: Recursos/2025-09-30.pdf, páginas 1--12.
En la clase 12 construimos productos y anillos de endomorfismos. Los ideales permiten multiplicar clases en un cociente y trasladar los teoremas de isomorfismo a los anillos. Estudiamos también generadores de ideales y cocientes por polinomios mónicos.

%%%%%%%%%%%%%%%%%%%%%%%%%%%%%%%%%%%%%%%%%%%%%%%%%%%%%%%%%%%%%%%%%%%%%%%%%%%%%%%
\section{Productos y endomorfismos}
%%%%%%%%%%%%%%%%%%%%%%%%%%%%%%%%%%%%%%%%%%%%%%%%%%%%%%%%%%%%%%%%%%%%%%%%%%%%%%%
%%%%%%%%%%%%%%%%%%%%%%%%%%%%%%%%%%%%%%%%%%%%%%%%%%%%%%%%%%%%%%%%%%%%%%%%%%%%%%%
\subsection{El producto en la categoría de anillos}
%%%%%%%%%%%%%%%%%%%%%%%%%%%%%%%%%%%%%%%%%%%%%%%%%%%%%%%%%%%%%%%%%%%%%%%%%%%%%%%
% P01, apartado III.2.4 y comprobación III.2.13.
Para anillos $R_1,R_2$, el conjunto $R_1\times R_2$ es anillo con operaciones por componentes:
\[
\begin{aligned}
(a_1,a_2)+(b_1,b_2)&=(a_1+b_1,a_2+b_2),\\
(a_1,a_2)(b_1,b_2)&=(a_1b_1,a_2b_2).
\end{aligned}
\]
Los axiomas se verifican en cada coordenada. Su cero es $(0,0)$ y su unidad es $(1_{R_1},1_{R_2})$.

\begin{ejer}[Propiedad universal del producto]\label{ejer:C12-producto}
Comprobar que este anillo, con sus proyecciones, es el producto categórico de $R_1$ y $R_2$ en $\Ring$.
\end{ejer}
\begin{sol}
Cada proyección $p_i$ preserva suma, producto y unidad, por lo que es morfismo. Dados homomorfismos $f_i:T\to R_i$, la única función con $p_i f=f_i$ es
\[
f:T\longrightarrow R_1\times R_2,\qquad t\longmapsto(f_1(t),f_2(t)).
\]
Preserva las operaciones por la definición componente a componente; también $f(1_T)=(1_{R_1},1_{R_2})$. Esto prueba la existencia y unicidad del mediador en $\Ring$:
\[
\begin{tikzcd}[column sep=large]
&T\arrow[dl,"f_1"']\arrow[d,dashed,destacado,"\exists!\,f"]\arrow[dr,"f_2"]&\\
R_1&R_1\times R_2\arrow[l,"p_1"]\arrow[r,"p_2"']&R_2.
\end{tikzcd}
\qedhere
\]
\end{sol}

%%%%%%%%%%%%%%%%%%%%%%%%%%%%%%%%%%%%%%%%%%%%%%%%%%%%%%%%%%%%%%%%%%%%%%%%%%%%%%%
\subsection{Anillos de endomorfismos y representación regular}
%%%%%%%%%%%%%%%%%%%%%%%%%%%%%%%%%%%%%%%%%%%%%%%%%%%%%%%%%%%%%%%%%%%%%%%%%%%%%%%
% P01--P03.
Sea $G$ un grupo abeliano. En $\End_{\Ab}(G)=\Hom_{\Ab}(G,G)$ definimos
\[
(\alpha+\beta)(g)=\alpha(g)+\beta(g),\qquad
\alpha\beta=\alpha\circ\beta.
\]
La suma de homomorfismos sigue siendo homomorfismo porque $G$ es abeliano. La suma puntual hace de este conjunto un grupo abeliano, con cero la función constante cero y negativo puntual. La composición es asociativa y tiene unidad $\id_G$. Las dos distributivas se comprueban evaluando: una usa la suma puntual y la otra la aditividad del endomorfismo que se compone a la izquierda. Por tanto, $\End_{\Ab}(G)$ es un anillo, en general no conmutativo.

% Mantener juntos el enunciado, su fórmula y el inicio de la prueba.
\begin{samepage}
\begin{prop}\label{prop:C12-endZ}
La evaluación en $1$ define un isomorfismo de anillos
\[
\Phi:\End_{\Ab}(\mathbb Z)\longrightarrow\mathbb Z,
\qquad\alpha\longmapsto\alpha(1).
\]
\end{prop}
\begin{proof}
Todo endomorfismo aditivo satisface $\alpha(n)=n\alpha(1)$. La evaluación preserva la suma y envía $\id_{\mathbb Z}$ a $1$. Además,
\[
\Phi(\alpha\circ\beta)=\alpha(\beta(1))
=\alpha(1)\beta(1)=\Phi(\alpha)\Phi(\beta).
\]
Su inversa es $\Psi(a)(n)=an$. En efecto, $\Psi(a)$ es aditiva, $\Psi(a+b)=\Psi(a)+\Psi(b)$, $\Psi(ab)=\Psi(a)\circ\Psi(b)$ y $\Psi(1)=\id$. Finalmente, $\Phi\Psi(a)=a$ y $\Psi\Phi(\alpha)(n)=\alpha(1)n=\alpha(n)$.
\end{proof}
\end{samepage}

Para un anillo $R$, las multiplicaciones izquierda y derecha son endomorfismos de $(R,+)$:
\[
\lambda_r(a)=ra,\qquad \mu_r(a)=ar.
\]
\begin{prop}\label{prop:C12-regular}
La regla $r\mapsto\lambda_r$ define un homomorfismo inyectivo de anillos $\lambda:R\to\End_{\Ab}(R)$.
\end{prop}
\begin{proof}
Por distributividad, $\lambda_{r+s}(a)=ra+sa=(\lambda_r+\lambda_s)(a)$. Por asociatividad,
\[
\lambda_{rs}(a)=(rs)a=r(sa)=(\lambda_r\circ\lambda_s)(a).
\]
También $\lambda_1=\id_R$. Finalmente, $\lambda_r(1)=r$, de modo que $\lambda_r=\lambda_s$ implica $r=s$.
\end{proof}
Esta representación es el análogo para anillos del argumento de Cayley. El orden de composición importa: para la multiplicación derecha se tiene $\mu_{rs}=\mu_s\circ\mu_r$, de modo que no se obtiene la misma afirmación sin invertir el producto.

%%%%%%%%%%%%%%%%%%%%%%%%%%%%%%%%%%%%%%%%%%%%%%%%%%%%%%%%%%%%%%%%%%%%%%%%%%%%%%%
\section{Ideales y cocientes}
%%%%%%%%%%%%%%%%%%%%%%%%%%%%%%%%%%%%%%%%%%%%%%%%%%%%%%%%%%%%%%%%%%%%%%%%%%%%%%%
%%%%%%%%%%%%%%%%%%%%%%%%%%%%%%%%%%%%%%%%%%%%%%%%%%%%%%%%%%%%%%%%%%%%%%%%%%%%%%%
\subsection{Absorción e imágenes inversas}
%%%%%%%%%%%%%%%%%%%%%%%%%%%%%%%%%%%%%%%%%%%%%%%%%%%%%%%%%%%%%%%%%%%%%%%%%%%%%%%
% P03.
\begin{defi}
Un subgrupo aditivo $I\subseteq R$ es \emph{ideal izquierdo} si $ra\in I$ para todo $r\in R$, $a\in I$; es \emph{ideal derecho} si $ar\in I$ para todos esos elementos. Es \emph{ideal bilateral} si cumple ambas condiciones. Cuando se diga simplemente ideal se entenderá bilateral.
\end{defi}
Los ideales están cerrados bajo sus propios productos, pero no se les exige contener $1_R$. Si un ideal izquierdo contiene $1_R$, entonces $r=r1_R\in I$ para todo $r$; para un ideal derecho se usa $r=1_Rr$. En ambos casos $I=R$. Por ello un ideal propio no es un subanillo unitario de $R$.

\begin{ejem}
El núcleo de cualquier homomorfismo $\varphi:R\to S$ es un ideal bilateral. Ya es subgrupo aditivo y, para $a\in\ker\varphi$, $r\in R$,
\[
\varphi(ra)=\varphi(r)\varphi(a)=\varphi(r)0=0,
\qquad
\varphi(ar)=0\varphi(r)=0.
\]
En particular, $\{0\}$ y $R$ son ideales de $R$.
\end{ejem}

\begin{ejer}[Preimagen de un ideal]\label{ejer:C12-preimagen}
Si $\varphi:R\to S$ es homomorfismo y $J$ es ideal de $S$, probar que $\varphi^{-1}(J)$ es ideal de $R$.
\end{ejer}
\begin{sol}
La preimagen contiene cero y, si $a,b$ pertenecen a ella, $\varphi(a-b)=\varphi(a)-\varphi(b)\in J$. Así es subgrupo aditivo. Para $r\in R$ y $a\in\varphi^{-1}(J)$,
\[
\varphi(ra)=\varphi(r)\varphi(a)\in J,\qquad
\varphi(ar)=\varphi(a)\varphi(r)\in J.
\]
Esto da las dos absorciones. La misma prueba conserva sólo el lado correspondiente si $J$ es unilateral. Tomando $J=\{0_S\}$ se recupera el ejemplo del núcleo.
\end{sol}

%%%%%%%%%%%%%%%%%%%%%%%%%%%%%%%%%%%%%%%%%%%%%%%%%%%%%%%%%%%%%%%%%%%%%%%%%%%%%%%
\subsection{Multiplicación de clases y característica}
%%%%%%%%%%%%%%%%%%%%%%%%%%%%%%%%%%%%%%%%%%%%%%%%%%%%%%%%%%%%%%%%%%%%%%%%%%%%%%%
% P04--P05.
Para un subgrupo aditivo $I$ de $R$, existe el grupo cociente $(R,+)/I$, porque el grupo es abeliano. Sus elementos son las clases $r+I$, con proyección aditiva $\pi:r\mapsto r+I$.

\begin{prop}\label{prop:C12-cociente}
La regla $(r+I)(s+I)=rs+I$ está bien definida si y sólo si $I$ es ideal bilateral. En ese caso hace de $R/I$ un anillo con unidad $1_R+I$ y convierte a $\pi:R\to R/I$ en un homomorfismo sobreyectivo de núcleo $I$.
\end{prop}
\begin{proof}
Supongamos que $I$ es ideal. Si $r-r'\in I$ y $s-s'\in I$, entonces
\[
rs-r's'=r(s-s')+(r-r')s'\in I.
\]
Por tanto, la clase del producto no depende de los representantes. La asociatividad, las distributivas y la identidad se heredan de $R$, y la proyección preserva ambas operaciones y la unidad.

Recíprocamente, si la regla está bien definida y $a\in I$, entonces $a+I=0+I$. Multiplicar esta igualdad por $r+I$ a cada lado da $ra+I=I=ar+I$. Así, $ra,ar\in I$ y $I$ es bilateral. Finalmente, $\pi(r)=I$ si y sólo si $r\in I$.
\end{proof}
\begin{obs}
El grupo aditivo de $R/I$ es un grupo cociente. No se identifica en general con un subgrupo o un subanillo de $R$.
\end{obs}

\begin{defi}
Sea $\iota_R:\mathbb Z\to R$ el morfismo inicial. Su núcleo es $n\mathbb Z$ para un único entero $n\geq0$. Se llama \emph{característica} de $R$ a ese entero, denotado $\operatorname{char}R$.
\end{defi}
Si $1_R$ tiene orden aditivo finito, la característica es ese orden: es el menor entero positivo $n$ con $n1_R=0$. Si tiene orden infinito, la característica es cero. En particular, el anillo cero tiene característica uno, $\mathbb Z$ tiene característica cero y $\mathbb Z/n\mathbb Z$ tiene característica $n$ para $n\geq1$.
\begin{obs}
El núcleo de $\iota_R$ puede ser trivial. Por eso el generador de ese subgrupo de $\mathbb Z$ es un entero $n\geq0$; el caso $n=0$ corresponde a característica cero.
\end{obs}

%%%%%%%%%%%%%%%%%%%%%%%%%%%%%%%%%%%%%%%%%%%%%%%%%%%%%%%%%%%%%%%%%%%%%%%%%%%%%%%
\subsection{Propiedad universal e isomorfismos}
%%%%%%%%%%%%%%%%%%%%%%%%%%%%%%%%%%%%%%%%%%%%%%%%%%%%%%%%%%%%%%%%%%%%%%%%%%%%%%%
% P06--P07.
\begin{teo}\label{teo:C12-universal}
Sea $I$ un ideal de $R$. Si $\varphi:R\to S$ es homomorfismo y $I\subseteq\ker\varphi$, existe un único homomorfismo $\widetilde\varphi:R/I\to S$ con $\widetilde\varphi\pi=\varphi$:
\[
\begin{tikzcd}[column sep=large,row sep=large]
R\arrow[r,"\varphi"]\arrow[dr,"\pi"']&S\\
&R/I\arrow[u,dashed,destacado,"\exists!\,\widetilde\varphi"'].
\end{tikzcd}
\]
\end{teo}
\begin{proof}
La propiedad universal del cociente aditivo proporciona la única aplicación aditiva $\widetilde\varphi(r+I)=\varphi(r)$. También se comprueba directamente: $r-r'\in I$ implica $\varphi(r)=\varphi(r')$. Para el producto,
\[
\widetilde\varphi((r+I)(s+I))=\varphi(rs)
=\varphi(r)\varphi(s)
=\widetilde\varphi(r+I)\widetilde\varphi(s+I).
\]
Además, $\widetilde\varphi(1_R+I)=\varphi(1_R)=1_S$. La conmutatividad fuerza su valor sobre cada clase, de modo que la unicidad aditiva también es unicidad como morfismo de anillos.
\end{proof}

\begin{teo}[Primer teorema de isomorfismo]\label{teo:C12-primer-iso}
Todo homomorfismo de anillos $\varphi:R\to S$ se factoriza como
\[
\begin{tikzcd}[column sep=large]
R\arrow[r,two heads,"\pi"]\arrow[rrr,bend left=24,"\varphi"]
&R/\ker\varphi\arrow[r,destacado,"\widetilde\varphi","\sim"']
&\im\varphi\arrow[r,hook]&S,
\end{tikzcd}
\]
donde $\widetilde\varphi(r+\ker\varphi)=\varphi(r)$. En particular, si $\varphi$ es sobreyectiva, $R/\ker\varphi\cong S$.
\end{teo}
\begin{proof}
La imagen es subanillo: es subgrupo aditivo, está cerrada bajo productos y contiene $1_S$. El \cref{teo:C12-universal} da el morfismo indicado con ese destino, y es sobreyectivo por definición de imagen. Su núcleo es trivial, porque $\varphi(r)=0$ equivale a $r\in\ker\varphi$. Por la \cref{prop:C11-mono} es inyectivo. La inversa de un homomorfismo biyectivo preserva las operaciones y la unidad al tomar preimágenes, así que es un isomorfismo de anillos.
\end{proof}

\begin{prop}[Cocientes sucesivos]\label{prop:C12-tercer-iso}
Si $I\subseteq J$ son ideales de $R$, entonces $J/I$ es ideal de $R/I$ y
\[
(R/I)/(J/I)\cong R/J.
\]
\end{prop}
\begin{proof}
La proyección $R\to R/J$ anula $I$, por lo que induce el homomorfismo sobreyectivo $q:R/I\to R/J$, $r+I\mapsto r+J$. Su núcleo es
\[
\ker q=\set{r+I\given r\in J}=J/I.
\]
Es un ideal por ser núcleo, y el \cref{teo:C12-primer-iso} produce el isomorfismo. El triángulo de factorización es el del \cref{teo:C12-universal}, con $S=R/J$.
\end{proof}

\begin{lema}[Correspondencia de ideales]\label{lema:C12-correspondencia}
Los ideales de $R/I$ corresponden, preservando inclusiones, a los ideales de $R$ que contienen $I$, mediante $J\mapsto J/I$ y $K\mapsto\pi^{-1}(K)$.
\end{lema}
\begin{proof}
Este detalle se usará en el criterio de maximalidad de la clase 13. La preimagen es ideal por el \cref{ejer:C12-preimagen} y contiene $I=\ker\pi$. Para $J\supseteq I$, la imagen $\pi(J)$ es $J/I$ y es ideal por la proposición anterior. Se tiene $\pi^{-1}(J/I)=J$: si $r+I=j+I$ con $j\in J$, entonces $r-j\in I\subseteq J$. Por la sobreyectividad de $\pi$, también $\pi(\pi^{-1}(K))=K$. Las dos operaciones respetan inclusiones, y son inversas.
\end{proof}

%%%%%%%%%%%%%%%%%%%%%%%%%%%%%%%%%%%%%%%%%%%%%%%%%%%%%%%%%%%%%%%%%%%%%%%%%%%%%%%
\section{Generadores de ideales y división de polinomios}
%%%%%%%%%%%%%%%%%%%%%%%%%%%%%%%%%%%%%%%%%%%%%%%%%%%%%%%%%%%%%%%%%%%%%%%%%%%%%%%
%%%%%%%%%%%%%%%%%%%%%%%%%%%%%%%%%%%%%%%%%%%%%%%%%%%%%%%%%%%%%%%%%%%%%%%%%%%%%%%
\subsection{Sumas, intersecciones y productos}
%%%%%%%%%%%%%%%%%%%%%%%%%%%%%%%%%%%%%%%%%%%%%%%%%%%%%%%%%%%%%%%%%%%%%%%%%%%%%%%
% P07--P10.
Para $a\in R$, $Ra=\set{ra\given r\in R}$ es ideal izquierdo y $aR$ es ideal derecho: son subgrupos aditivos, y la absorción se sigue de $s(ra)=(sr)a$ y $(ar)s=a(rs)$. Si $R$ es conmutativo, coinciden y se denotan $(a)$, el \emph{ideal principal generado por $a$}. En particular, $(0)=\{0\}$ y $(1)=R$.

\begin{ejer}[Suma de ideales]\label{ejer:C12-suma}
Para una familia $(I_\alpha)_{\alpha\in A}$ de ideales de un anillo $R$, probar que su suma es un ideal.
\end{ejer}
\begin{sol}
Definimos la suma mediante \emph{sumas finitas} de elementos de los ideales:
\[
\sum_{\alpha\in A}I_\alpha
=\set*{a_1+\cdots+a_k\given
 k\geq0,\ a_i\in I_{\alpha_i},\ \alpha_i\in A}.
\]
La suma vacía es cero. Sumar dos expresiones concatena las listas, y negar una expresión niega cada sumando. Así es subgrupo aditivo. Multiplicar por cualquier $r\in R$ a la izquierda o a la derecha deja cada sumando en su ideal original, por lo que la suma es bilateral. Contiene a cada $I_\alpha$, y cualquier ideal que los contenga incluye todas esas sumas finitas. Es, por tanto, el menor ideal con esa propiedad. Para $A$ vacío se obtiene $(0)$.
\end{sol}

Si $R$ es conmutativo y $(a_\alpha)_{\alpha\in A}$ es una familia de elementos, definimos el \emph{ideal generado} por ellos como $\sum_{\alpha\in A}(a_\alpha)$. Para una familia finita,
\[
(a_1,\ldots,a_n)
=\set{r_1a_1+\cdots+r_na_n\given r_1,\ldots,r_n\in R}.
\]
Es el menor ideal que contiene todos los generadores. Un ideal es \emph{finitamente generado} si admite una expresión de este tipo.

\begin{ejem}[Añadir una relación después de formar un cociente]
Para $R$ conmutativo, $a,b\in R$ y $\bar b=b+(a)$, se tiene
\[
(R/(a))/(\bar b)\cong R/(a,b).
\]
En efecto, en $R/(a)$,
\[
(\bar b)=\set{zb+(a)\given z\in R}=(a,b)/(a).
\]
La última igualdad se debe a que $ua+vb+(a)=vb+(a)$: añadir un múltiplo de $a$ no cambia la clase. Aplicamos la \cref{prop:C12-tercer-iso} a $(a)\subseteq(a,b)$.
\end{ejem}

La intersección $\bigcap_{\alpha\in A}I_\alpha$ es el mayor ideal contenido en cada $I_\alpha$: las condiciones de subgrupo y absorción se comprueban en todos ellos. Si la familia es vacía, su intersección es $R$.

El \emph{producto de ideales} $IJ$ es el ideal generado por los productos $ab$ con $a\in I$, $b\in J$. Para ideales bilaterales de un anillo cualquiera consiste en las sumas finitas de tales productos: al multiplicarlas por un elemento de $R$, la absorción se aplica al primer o al segundo factor. Más generalmente, $I_1\cdots I_n$ es generado por los productos $a_1\cdots a_n$, $a_k\in I_k$. Esta notación no designa sólo el conjunto de productos individuales, a diferencia del producto de subgrupos empleado antes.

%%%%%%%%%%%%%%%%%%%%%%%%%%%%%%%%%%%%%%%%%%%%%%%%%%%%%%%%%%%%%%%%%%%%%%%%%%%%%%%
\subsection{Anillos noetherianos y dominios de ideales principales}
%%%%%%%%%%%%%%%%%%%%%%%%%%%%%%%%%%%%%%%%%%%%%%%%%%%%%%%%%%%%%%%%%%%%%%%%%%%%%%%
% P09.
\begin{defi}
Un anillo conmutativo es \emph{noetheriano} si todos sus ideales son finitamente generados. Un \emph{dominio de ideales principales}, abreviado DIP, es un dominio entero cuyos ideales son todos principales.
\end{defi}
Todo DIP es noetheriano, ya que un generador es una familia finita.

\begin{prop}\label{prop:C12-ZDIP}
El anillo $\mathbb Z$ es un dominio de ideales principales.
\end{prop}
\begin{proof}
Es dominio. Cada ideal es subgrupo aditivo de $\mathbb Z$, luego es $n\mathbb Z$ para algún $n\geq0$, por la clasificación de subgrupos de un grupo cíclico. Pero $n\mathbb Z=(n)$, así que es principal.
\end{proof}

\begin{ejem}[Máximo común divisor]
En $\mathbb Z$, $(m,n)=(d)$ con $d\geq0$. Como $m,n\in(d)$, $d$ divide a ambos. Como $d\in(m,n)$, existen $a,b\in\mathbb Z$ con $d=am+bn$. Todo divisor común $d'$ divide esa combinación, luego divide a $d$. Por tanto, $d=\mcd(m,n)$ y la descripción del ideal produce la identidad de Bézout. Si $m=n=0$, se toma $d=0$.
\end{ejem}

\begin{ejem}\label{ejem:C12-noDIP}
El anillo $\mathbb Z[x]$ no es un DIP: el ideal $(2,x)$ no es principal. Primero, $(2,x)$ es propio, porque el término constante de $2a(x)+xb(x)$ es par y, por tanto, $1$ no pertenece al ideal.

Si $(2,x)=(f)$, entonces $f$ divide a $x$ y a $2$. En $\mathbb Z[x]$ el grado de un producto no nulo es la suma de los grados, porque el producto de sus coeficientes principales no se anula. La igualdad $x=fg$ obliga a que $f=\pm1$ o $f=\pm x$: uno de los factores es constante y la igualdad de coeficientes principales exige constantes unidades. La primera posibilidad contradice que el ideal sea propio; la segunda contradice que $2\in(f)$, pues un múltiplo de $x$ tiene término constante cero.
\end{ejem}
\begin{obs}
En $(2,x)\subseteq\mathbb Z[x]$, los coeficientes de una combinación $2a+bx$ pertenecen a $\mathbb Z[x]$. Por ejemplo, $x^2$ pertenece al ideal.
\end{obs}

%%%%%%%%%%%%%%%%%%%%%%%%%%%%%%%%%%%%%%%%%%%%%%%%%%%%%%%%%%%%%%%%%%%%%%%%%%%%%%%
\subsection{División por un polinomio mónico}
%%%%%%%%%%%%%%%%%%%%%%%%%%%%%%%%%%%%%%%%%%%%%%%%%%%%%%%%%%%%%%%%%%%%%%%%%%%%%%%
% P10--P12.
En este apartado $R$ es conmutativo y no nulo. El grado de un polinomio no nulo es el mayor exponente con coeficiente no nulo; convenimos que $\deg0=-\infty$. Un polinomio es \emph{mónico} si su coeficiente principal es $1_R$.

\begin{lema}\label{lema:C12-monico}
Si $f$ es mónico de grado $d$ y $h\ne0$, entonces $fh\ne0$ y $\deg(fh)=d+\deg h$. En particular, $f$ no es divisor de cero. Si
\[
fq_1+r_1=fq_2+r_2,\qquad \deg r_1,\deg r_2<d,
\]
entonces $q_1=q_2$ y $r_1=r_2$.
\end{lema}
\begin{proof}
Si el término principal de $h$ es $cx^m$, con $c\ne0$, el coeficiente de $x^{d+m}$ en $fh$ es $1_Rc=c\ne0$. Esto demuestra la primera afirmación incluso si $R$ tiene divisores de cero.

Para la segunda, $f(q_1-q_2)=r_2-r_1$. Si $q_1-q_2\ne0$, el lado izquierdo tiene grado al menos $d$, mientras que el derecho tiene grado menor que $d$ o es cero. La contradicción da $q_1=q_2$, de donde se sigue $r_1=r_2$.
\end{proof}

% El enunciado no debe quedar aislado al pie de la página anterior.
\begin{samepage}
\begin{prop}[División por un mónico]\label{prop:C12-division}
Para $g\in R[x]$ y $f$ mónico de grado $d$, existen únicos $q,r\in R[x]$ tales que $g=fq+r$ y $\deg r<d$.
\end{prop}
\begin{proof}
La unicidad es el \cref{lema:C12-monico}. Aplicamos el algoritmo de división. Si $g=0$ o $\deg g<d$, tomamos $q=0$, $r=g$. Si el término principal de $g$ es $cx^m$ con $m\geq d$, restamos $cx^{m-d}f$: el término de grado $m$ desaparece, por lo que el nuevo polinomio tiene grado menor o es cero. Repetimos el procedimiento. El grado desciende estrictamente en cada paso no trivial, de modo que termina y las cantidades restadas reúnen el cociente $q$. Para $d=0$, $f=1$ y la única posibilidad es $q=g$, $r=0$.
\end{proof}
\end{samepage}

\begin{prop}\label{prop:C12-restos}
Si $f\in R[x]$ es mónico de grado $d$, cada clase de $R[x]/(f)$ tiene un único representante de grado menor que $d$. Sus coeficientes inducen un isomorfismo de grupos abelianos
\[
R[x]/(f)\cong R^{\oplus d}.
\]
\end{prop}
\begin{proof}
La división $g=fq+r$ da $g+(f)=r+(f)$. Si dos restos representan la misma clase, su diferencia es múltiplo de $f$ y tiene grado menor que $d$; el \cref{lema:C12-monico} obliga a que sea cero.

Identificamos $r=a_0+a_1x+\cdots+a_{d-1}x^{d-1}$ con $(a_0,\ldots,a_{d-1})$. La aplicación que asigna a $g$ ese vector es aditiva: si $g_i=fq_i+r_i$, entonces $g_1+g_2=f(q_1+q_2)+(r_1+r_2)$, y la suma de restos sigue teniendo grado menor que $d$. Es sobreyectiva, porque todo vector determina un resto. Su núcleo es $(f)$, pues el resto es cero exactamente para los múltiplos de $f$. El primer teorema de isomorfismo de \emph{grupos aditivos} da el resultado. Para $d=0$, ambos lados son el grupo cero.
\end{proof}
Este isomorfismo describe la suma. Para multiplicar clases se multiplican sus representantes y se toma el resto del producto; en general, esto difiere del producto por componentes.

\begin{obs}[Referencias de lectura]
Para productos, ideales, cocientes y polinomios, véase \citet[cap.~III, 2.4--7, 3.1/3/8--11 y 4.1--6]{aluffi2009}. Los ejercicios corresponden a III.2.13, III.3.2 y III.4.1.
\end{obs}
